\documentclass[aspectratio=169]{beamer}
\usepackage{tikz}
\usepackage{media9}
\usepackage{color}
\definecolor{theme}{RGB}{0,73,114}

\newcommand{\tabitem}{%
	\usebeamertemplate{itemize item}\hspace*{\labelsep}}

\usetikzlibrary{backgrounds,
	calc,
	fadings,
	shadows,
	shapes.arrows,
	shapes.symbols
}

\usepackage{xparse}
\edef\RestoreEndlinechar{%
	\endlinechar=\the\endlinechar\relax
}
\endlinechar=-1
\usepackage{smartdiagram}
\RestoreEndlinechar

\tikzset{bubble node/.append style={
		draw=none, opacity=0.5
	}
}

\definecolor{BlueCLS}{HTML}{00AAFF}
\definecolor{DarkBlueCLS}{HTML}{122372}
\definecolor{GrayCLS}{HTML}{B1B3B4}
\definecolor{LightGrayCLS}{HTML}{D8D9D9}
\definecolor{YellowCLS}{HTML}{FFED00}

\newcommand{\Aback}{0,0}
\newcommand{\Bback}{{2.9*\paperwidth/29.7},0}
\newcommand{\Cback}{{3.6*\paperwidth/29.7},{3*\paperheight/21}}
\newcommand{\Dback}{0,{1.9*\paperheight/21}}
\newcommand{\Eback}{{21.2*\paperwidth/29.7},{8.5*\paperheight/21}}
\newcommand{\Fback}{{7.5*\paperwidth/29.7},\paperheight}
\newcommand{\Gback}{0,{14.1*\paperheight/21}}
\newcommand{\Hback}{0,\paperheight}
\newcommand{\Iback}{\paperwidth,{1*\paperheight/21}}
\newcommand{\Jback}{\paperwidth,{11*\paperheight/21}}
\newcommand{\Kback}{\paperwidth,0}
\newcommand{\Scale}{0.01*\paperwidth/29.7}

\setbeamertemplate{navigation symbols}{}
\graphicspath{{./images/}}

\begin{document}
	
	% Title slide frame
\begin{frame}[plain]	
		%%%%%%%% Title slide details %%%%%%%%%%%%%%


% Background Image
\newcommand{\myBackground}
{
\begin{tikzpicture}[anchor=south east, yscale=-1]
 	% Yellow ABCD
 	\coordinate (A) at (\Aback);
 	\coordinate (B) at (\Bback);
 	\coordinate (C) at (\Cback);
 	\coordinate (D) at (\Dback);
 	% Gray DEFG
 	\coordinate (E) at (\Eback);
 	\coordinate (F) at (\Fback);
 	\coordinate (G) at (\Gback);
 	% DarkBlue GFH
 	\coordinate (H) at (\Hback);
 	% Blue EIJ
 	\coordinate (I) at (\Iback);
 	\coordinate (J) at (\Jback);
 	% LightGray BKIEC
 	\coordinate (K) at (\Kback);
 	% Background
 	\fill[YellowCLS, opacity=0.4] (A) -- (B) -- (C) -- (D) -- cycle;
 	\fill[GrayCLS, opacity=0.4] (D) -- (E) -- (F) -- (G) -- cycle;
 	\fill[DarkBlueCLS, opacity=0.4] (G) -- (F) -- (H) -- cycle;
 	\fill[BlueCLS, opacity=0.4] (E) -- (I) -- (J) -- cycle;
 	\fill[LightGrayCLS, opacity=0.4] (B) -- (K) -- (I) -- (E) -- (C) -- cycle;
 	%\raisebox{-0.79\paperheight}{\hspace{0.82\paperwidth}~\insertLogo}
\end{tikzpicture}
}

% Title
\newcommand{\myTitle}
{
   Une formulation port-Hamiltonienne des structures flexibles. \\
   Modélisation et discrétisation symplectique par éléments finis
}

% Subtitle
\newcommand{\mySubTitle}
{

}

% Author
\newcommand{\myAuthor}   
{
    Andrea Brugnoli
}

% Affiliation
\newcommand{\myAffiliate}
{
  Sous la direction de: D. Alazard, V. Budinger (DCAS) et D. Matignon (DISC)
}

% Presentation Date
\newcommand{\myAck}   
{
 Dessins: F. Sanfedino
}

% Logo
\newcommand{\myLogo}   
{
	\begin{tikzpicture}[remember picture, overlay]
		\node[anchor=south, shift={(0\paperwidth,0.03\paperheight)}] at (current page.south) {\includegraphics[width=3cm]{ISAE-SUPAERO.png}
		};	
	\end{tikzpicture}


}
%%%%%%%%%%%%%%%%%%%%%%%%%%%%%%%%%%%%


%%%%%%%%%% Title slide code %%%%%%%%%%%
\begin{tikzpicture}[remember picture,overlay]

% Background color

\fill[white] (current page.south west) rectangle (current page.north east);
% Background image
\node[above right,inner sep=0pt] at (current page.south west)
    {
        \myBackground
    };
    
% Title & Subtitle
\node
[
    above=0.5cm,
    align=center,
    draw=black!50,
    % rounded corners,
    double,
    double distance=0.1cm,
    double=blue!10,
    fill=theme!10,
    inner xsep=15pt,
    inner ysep=10pt, 
    minimum width=0.7\textwidth,
    text width=0.9\textwidth
] (title) at (current page.center)
{
    \Large \myTitle  \\[5pt]
    \small \mySubTitle
};

% Author 
\node[ below=0.5cm] (author) at (title.south){\Large \myAuthor};

% Affiliate 
\node[below=0.25cm](affiliate) at (author.south){\small \myAffiliate};

% Ack
\node[below=0.2cm] (date) at (affiliate.south){\small \myAck};

% Logo
\node
[
    below =0.25cm
] at (date.south)
{
    \myLogo
};

\end{tikzpicture}
\end{frame}

{
	\setbeamercolor{background canvas}{bg=black}
	\begin{frame}<article:0>[plain]
		\begin{tikzpicture}[remember picture,overlay]
			\node[at=(current page.center)] {
				\includegraphics[keepaspectratio,
				width=\paperwidth,
				height=\paperheight]{satellite_panels.jpg}
			};
		\end{tikzpicture}
	\end{frame}
}


{
	\setbeamercolor{background canvas}{bg=black}
	\begin{frame}<article:0>[plain]
		\begin{tikzpicture}[remember picture,overlay]
			\node[at=(current page.center)] {
				\includegraphics[keepaspectratio,
				width=\paperwidth,
				height=\paperheight]{helios.jpg}
			};
		\end{tikzpicture}
	\end{frame}
}


\begin{frame}<article:0>[plain]
	\begin{tikzpicture}[remember picture,overlay]
		\node[at=(current page.center)] {
			\includemedia[
			label=vidNoRod,
			addresource=/home/andrea/Videos/PrixThese/flutter_final.mp4,
			activate=pageopen,
			keepaspectratio,  width=\paperwidth, height=\paperheight,
			flashvars={
				source=/home/andrea/Videos/PrixThese/flutter_final.mp4
				&loop=true
			}
			]{}{VPlayer.swf}
		};
	\end{tikzpicture}
\end{frame}

\begin{frame}

\begin{center}
	\begin{tabular}{@{}l@{}}
		Outil de modélisation: \\
		\tabitem intuitif \\
		\tabitem respectueux de la physique
	\end{tabular}
\end{center}

\end{frame}

\begin{frame}
	Immagine aereo modulare 
\end{frame}

\begin{frame}
	Immagine aereo e radar e vortici 
\end{frame}

\begin{frame}
	\begin{columns}
		\begin{column}{.5\textwidth}
			\smartdiagramset{
				distance text center bubble=0.5cm,
				bubble center node size=1cm,
				bubble node size=2cm,
				distance center/other bubbles=0.5cm,
				bubble center node font=\footnotesize,
				bubble node font=\scriptsize,
				bubble center node color=cyan,
				set color list = {yellow, pink, blue, red}
			}%
			\smartdiagram[bubble diagram]{%
				Systèmes \\
				Hamiltoniens.,
				\only<1>{Mecanique}\only<2>{\textcolor{red}{\large Mecanique}},
				Dynamique \\
				des fluides,
				Électromagnétisme,
				Thermodynamique%
			}%
		\end{column}
	\begin{column}{.5\textwidth}
		\only<2>{
		\includegraphics[width=\columnwidth]{satellite.pdf}
	}
	\end{column}
	\end{columns}
	
	
\end{frame}


\begin{frame}
	\only<1>{
	\includegraphics[width=0.5\textwidth]{cube.pdf}\\
	Domaine physique}
	\only<2>{
	\includegraphics[width=0.5\textwidth]{mesh_cube.pdf}\\
	Maillage}
\end{frame}

\end{document}